% =============================================================================
% cap01.tex — INTRODUCAO
% =============================================================================

\chapter{INTRODUÇÃO}\label{cap:introducao}
% 00-siglas.tex — siglas registradas na lista sem aparecer por extenso aqui (issue 98).
% As que aparecem por extenso no texto usam \criarsigla* no ponto da 1ª ocorrência
% (MLC, NDP, ABNT, CEN, NPT, CBMPR, CCA, CCB, BDI); EC5 é apresentada junto do CEN.
\criarsigla{EC5}{Eurocódigo 5, EN 1995-1-1: projeto de estruturas de madeira}
\criarsigla{EN}{norma europeia}
\criarsigla{NBR}{norma brasileira}
\criarsigla{ISO}{Organização Internacional de Normalização}
\criarsigla{UC}{classe de uso, EN 335}
\criarsigla{NP}{classe de penetração do preservativo, EN 351-1}
\criarsigla{CV-H}{contraventamento horizontal, da cobertura}
\criarsigla{CV-V}{contraventamento vertical, painéis das fachadas e do alinhamento central}
\criarsigla{CV-T}{contraventamento transversal, painéis dos frontões}
\criarsigla{CV-D}{alma do diafragma da cobertura, cruzes nas baías intermediárias}
\criarsigla{L1, L1c, L1f}{ligações de topo: viga sobre pilar lateral, sela do pilar central e topo do pilar de frontão}
\criarsigla{L3, L3c}{ligações de base correntes: pilar lateral e pilar central}
\criarsigla{L3f, L3a}{ligações de base do pilar de frontão e base ancorada dos montantes}
\criarsigla{P, PF, T, V}{pilares, pilares de frontão, terças e vigas, numerados conforme a tabela Nomenclatura dos elementos, que define também as designações acima}


A madeira voltou a ocupar lugar de destaque entre os materiais estruturais. Às
qualidades que sempre a recomendaram, como a elevada relação entre resistência e
peso, a facilidade de trabalho e a pré-fabricação, a montagem rápida e a seco e o
apelo estético das superfícies aparentes, somou-se nas últimas décadas uma
vantagem ambiental. A madeira provém de fonte renovável, demanda pouca energia
na produção quando comparada ao aço e ao concreto e mantém armazenado, durante
toda a vida útil da edificação, o carbono que a árvore retirou da atmosfera
\cite{ramage:2017}. Substituir por madeira parte dos materiais convencionais nas
novas construções é apontado como um meio de converter o estoque edificado em
reservatório de carbono \cite{churkina:2020}.

Esse movimento tem como protagonista a chamada madeira maciça engenheirada,
designação dos produtos industrializados que permitem empregar a madeira em
elementos de grandes dimensões, com propriedades controladas
\cite{harte:2017}. O mais antigo e difundido deles é a
\criarsigla*{MLC}{madeira lamelada colada}, que viabilizou, ao longo do último
século, coberturas e pórticos de vãos que a madeira serrada não alcança
\cite{herzog:2004}.

A MLC é um
produto estrutural obtido pela colagem de lamelas de madeira serrada com as
fibras orientadas paralelamente ao eixo longitudinal da peça. O processo de
laminação dissocia a dimensão da peça acabada da dimensão da tora de origem e
permite, ao mesmo tempo, a remoção seletiva dos defeitos das lamelas antes da
colagem. Disso decorrem as duas características que sustentam o emprego
estrutural do material: a possibilidade de produzir elementos de grande vão e
de seção variável, e uma dispersão de resistência menor do que a da madeira
serrada de mesma espécie, o que se traduz em valores característicos mais
elevados \cite{thelandersson:2003, porteous:2013}.

Essas características tornam a MLC competitiva justamente na faixa de
aplicação em que a madeira serrada não alcança e em que as soluções em aço ou
concreto se tornam onerosas: coberturas de grande vão livre, com carregamento
predominantemente permanente e sobrecarga reduzida. Galpões industriais,
pavilhões, centros de distribuição e edifícios-garagem constituem o campo
típico dessa aplicação.

%% TODO: verificar — referências inseridas na issue 57 (ramage:2017,
%% churkina:2020, harte:2017, porteous:2013, thelandersson:2003, herzog:2004):
%% conferir os dados bibliográficos com as fontes antes da entrega. Falta
%% literatura nacional sobre o mercado brasileiro de MLC.

\section{CONTEXTO NORMATIVO}\label{sec:contexto-normativo}

A primeira norma brasileira de estruturas de madeira, a NB-11 de 1951
\cite{NB11:1951}, tinha como
método de segurança o das tensões admissíveis. A passagem para o método dos
estados-limites, de natureza probabilística, ocorreu com a ABNT NBR 7190:1997
\cite{NBR7190:1997}, norma da
\criarsigla*{ABNT}{Associação Brasileira de Normas Técnicas} que já organizava o projeto em verificações de
estados-limites últimos e de serviço, com coeficientes parciais de segurança e
coeficientes de modificação da resistência. Não é esse, portanto, o traço que
distingue a revisão de 2022 \cite{NBR7190-1:2022}.

O que a revisão de 2022 trouxe foi a atualização dos critérios de dimensionamento
com base em ensaios de novos materiais e ligações. A norma passou a ser publicada
em sete partes: a ABNT NBR 7190-1 reúne os critérios de dimensionamento, e as
Partes 2 a 7 tratam dos métodos de ensaio, respectivamente para a classificação
visual e mecânica de peças estruturais, para corpos de prova isentos de defeitos
de madeiras nativas, para a caracterização de peças estruturais, para ligações
com conectores mecânicos, para a madeira lamelada colada e para a madeira
lamelada colada cruzada. As classes de resistência foram redefinidas a partir
de ensaios de peças estruturais, em série semelhante à europeia, com as classes
C14 a C50 para coníferas e D18 a D70 para folhosas, mantida para as espécies
nativas a caracterização por corpos de prova isentos de defeitos, nas classes
D20 a D60. O coeficiente de modificação passou a ser o produto de dois fatores,
$k_{mod} = k_{mod1} \cdot k_{mod2}$, o primeiro dependente da classe de
carregamento e do tipo de material, com valores explícitos para a madeira
lamelada colada, a lamelada colada cruzada e a laminada colada, e o segundo da
classe de umidade. Foram incluídos ainda novos coeficientes de modificação, um
novo critério para a estabilidade de peças comprimidas e flexocomprimidas e
recomendações para a segurança em situação de incêndio.

A norma brasileira, entretanto, não fornece o procedimento de verificação da
tipologia deste trabalho: vigas de MLC com dupla inclinação devem, segundo ela,
ser analisadas individualmente, sem as regras para a borda inclinada e para as
tensões transversais na zona de cumeeira que o Eurocódigo apresenta. As
diferenças entre as duas abordagens são discutidas na
\autoref{sec:eurocodigo-nbr}.

A EN 1995-1-1 \cite{EN1995-1-1:2004}, norma do
\criarsigla*{CEN}{Comitê Europeu de Normalização} conhecida como Eurocódigo 5
(EC5), cobre essas verificações
de maneira explícita e é complementada por um conjunto normativo coerente: a
EN 14080 \cite{EN14080:2013}, que define as classes de resistência e as propriedades
características da MLC; a EN 1990 \cite{EN1990:2002}, que estabelece as combinações
de ações; e as partes pertinentes da EN 1991 para os carregamentos. Por esse
motivo, o presente trabalho adota o Eurocódigo 5 como base de dimensionamento,
recorrendo à ABNT NBR 6123 \cite{NBR6123:2023} para a determinação das forças devidas ao
vento, por ser esta a norma que descreve as condições meteorológicas do sítio
brasileiro.

Essa composição normativa não é neutra e introduz uma dificuldade prática: as
verificações do Eurocódigo 5 dependem de
\criarsigla*{NDP}{parâmetros de determinação nacional}, que deveriam ser
fixados por um Anexo Nacional inexistente no Brasil. O tratamento dado a essa
lacuna é discutido adiante, no \autoref{cap:definicao-problema}.

\section{JUSTIFICATIVA}\label{sec:justificativa}

A verificação de um pórtico de MLC segundo o Eurocódigo 5 não se resume ao
dimensionamento da viga principal. Uma estrutura completa exige, além da flexão e
do cisalhamento da viga, a verificação das tensões na cumeeira e na borda
inclinada, da estabilidade lateral, dos pilares, das terças, da fixação da
telha sob sucção, do contraventamento horizontal e vertical, da estabilidade
global, das ligações e da ancoragem à fundação. Cada uma dessas verificações
depende de resultados das demais, e uma alteração em qualquer variável
geométrica propaga-se por todo o conjunto.

Executada manualmente, essa cadeia de verificações é lenta o bastante para
inviabilizar o estudo de alternativas: na prática, adota-se a primeira
geometria que atende a todos os critérios, sem que se saiba a que distância ela
está de uma solução eficiente. A automação do memorial de cálculo remove essa
restrição, pois permite varrer uma grade de geometrias candidatas, identificar
qual verificação governa em cada região dessa grade e submeter cada etapa a
testes automatizados que reduzem o risco de regressões nos comportamentos que
cobrem.

\section{OBJETIVOS}\label{sec:objetivos}\label{sec:objetivo-geral}

O objetivo deste trabalho é dimensionar e verificar a estrutura de uma cobertura
para estacionamento em madeira lamelada colada com pórticos de vigas de duplo-bisel, segundo a
EN 1995-1-1 \cite{EN1995-1-1:2004}, por meio de um programa de cálculo
automatizado e reprodutível, e apresentar o memorial resultante. São objetivos
específicos:\label{sec:objetivos-especificos}

\begin{alineas}
    \item obter a geometria da viga de duplo-bisel e as seções dos pilares e das
    terças que atendem ao conjunto de verificações de estado-limite último e de
    serviço;

    \item definir os sistemas de contraventamento que estabilizam a cobertura e a
    edificação e conduzem as forças laterais até a fundação;

    \item detalhar as ligações e a ancoragem, em desenhos cotados e em escala, e
    fornecer os esforços de cálculo destinados ao projeto da fundação;

    \item estimar o custo preliminar da estrutura;

    \item disponibilizar o programa de cálculo, com código aberto e testes
    automatizados, capaz de reproduzir o memorial a partir dos dados de entrada.
\end{alineas}

\section{DELIMITAÇÃO DO TRABALHO}\label{sec:delimitacao}

O trabalho delimita-se à superestrutura em madeira lamelada colada e aos seus
elementos de ligação. A fundação não é dimensionada: o programa entrega os
esforços característicos e de cálculo em cada base, a saber, reação vertical,
força horizontal e esforço de arrancamento, que constituem os dados de entrada
para o projeto geotécnico, executado em etapa posterior e a partir de sondagem
específica do terreno.

Também não integram o escopo o projeto de instalações, o detalhamento de
fechamentos e vedações e o plano de montagem.

A resistência ao fogo não é exigida da estrutura. O tempo requerido não é
fixado pelo Eurocódigo, que dá apenas o método de verificação. Ele decorre da
ABNT NBR 14432 \cite{NBR14432:2001} e, no Paraná, da \criarsigla*{NPT}{Norma de Procedimento Técnico} 008 do
\criarsigla*{CBMPR}{Corpo de Bombeiros Militar do Paraná}, na versão de 2026
\cite{CBMPR-NPT008:2026}. Por essa norma,
a cobertura, térrea e aberta nas quatro faces, destinada a estacionamento, é
isenta de tempo requerido de resistência ao fogo, conforme a
\autoref{ssec:trrf}. A verificação da seção reduzida a 30 minutos é mantida como
informativa. A mesma norma exige, independentemente da isenção, retardante ao fogo
nos elementos de madeira, especificado na \autoref{ssec:faces-expostas}.

\section{ESTRUTURA DO TRABALHO}\label{sec:estrutura-trabalho}

O texto está organizado em oito capítulos. Após esta introdução, o
\autoref{cap:revisao} reúne a revisão bibliográfica, e o
\autoref{cap:definicao-problema} apresenta a edificação, os dados de entrada e as
reservas normativas assumidas. A partir daí, a exposição acompanha a sequência
do programa de cálculo: o \autoref{cap:predimensionamento} trata da varredura
que conduziu à geometria adotada; o \autoref{cap:modelagem}, do modelo
estrutural, das ações e dos esforços; o \autoref{cap:dimensionamento}, das
verificações de estado-limite último e de serviço; e o
\autoref{cap:detalhamento}, das ligações, do contraventamento e da ancoragem. O
\autoref{cap:consideracoes} encerra o texto. Os apêndices descrevem o programa
(\autoref{ap:codigo-fonte}) e reúnem a memória de cálculo das verificações
governantes (\autoref{ap:memoria-calculo}), a análise estática da viga
(\autoref{ap:analise-estatica}), a saída do programa
(\autoref{ap:saida-pipeline}) e os desenhos (\autoref{ap:desenhos}).

Em todo o texto, dados de entrada e resultados ocupam tabelas separadas. Os
dados gerais concentram-se no \autoref{cap:definicao-problema}; os de um único
elemento, como o arranjo de uma ligação, aparecem no capítulo em que são usados,
antes da tabela de resultados correspondente.
